\documentclass[12pt]{article}
\usepackage{ficha}

\begin{document}

\ficha{Potências e Raízes}[Potências de Expoente Natural, Quadrados, Cubos e Raízes]

\instrucoes{Resolve sem calculadora, apresentando os cálculos. Apresenta os resultados na forma de uma única potência sempre que pedido.}

\secao{Potências de Expoente Natural}{Escreve na forma de potência ou calcula o valor, conforme pedido.}
\begin{exercicios}
  \item Escreve na forma de potência: $5 \times 5 \times 5 \times 5$.
  \item Escreve na forma de potência: $(-2) \times (-2) \times (-2)$.
  \item Indica a base e o expoente de $7^4$.
  \item Calcula $2^5$.
  \item Calcula $3^4$.
  \item Calcula $10^6$.
  \item Calcula $1^{20}$ e $0^{7}$.
  \item Calcula $(-3)^2$ e $(-3)^3$.
  \item Calcula $\left(\dfrac{2}{3}\right)^2$.
  \item Calcula $(0{,}5)^3$.
  \item Qual é maior: $2^5$ ou $5^2$?
\end{exercicios}

\secao{Quadrados e Cubos Perfeitos}{Responde a cada questão.}
\begin{exercicios}
  \item Escreve os quadrados perfeitos de $1^2$ a $15^2$.
  \item Escreve os cubos perfeitos de $1^3$ a $10^3$.
  \item 144 é um quadrado perfeito? De que número?
  \item 125 é um cubo perfeito? De que número?
  \item Indica um número que seja simultaneamente quadrado e cubo perfeito, além de 1.
  \item Um quadrado tem área 81\,cm$^2$. Quanto mede o lado?
  \item Um cubo tem volume 216\,cm$^3$. Quanto mede a aresta?
  \item Entre que dois quadrados perfeitos consecutivos está 70?
\end{exercicios}

\secao{Regras Operatórias}{Escreve cada expressão na forma de uma única potência.}
\begin{exercicios}
  \item $2^3 \times 2^4$
  \item $5^7 : 5^3$
  \item $(3^2)^4$
  \item $4^3 \times 5^3$
  \item $12^5 : 4^5$
  \item $7^6 \times 7 : 7^4$
  \item $(2^3)^2 \times 2^5$
  \item $(-3)^4 \times (-3)^5$
  \item $\dfrac{10^8}{10^5}$
  \item $6^4 : 2^4 \times 3^2$
  \item $(5^2 \times 5^3)^2 : 5^6$
  \item $\left(\dfrac{1}{2}\right)^3 \times \left(\dfrac{1}{2}\right)^4$
\end{exercicios}

\secao{Raízes Quadradas e Cúbicas}{Calcula.}
\begin{exercicios}
  \item $\sqrt{49}$
  \item $\sqrt{121}$
  \item $\sqrt{400}$
  \item $\sqrt[3]{27}$
  \item $\sqrt[3]{1\,000}$
  \item $\sqrt[3]{-64}$
  \item $\sqrt{\dfrac{25}{36}}$
  \item $\sqrt{0{,}09}$
  \item $\sqrt{9 + 16}$ \quad e \quad $\sqrt{9} + \sqrt{16}$. São iguais?
  \item Entre que dois números inteiros consecutivos está $\sqrt{50}$?
  \item Entre que dois números inteiros consecutivos está $\sqrt[3]{100}$?
\end{exercicios}

\secao{Expressões com Potências e Raízes}{Calcula o valor de cada expressão.}
\begin{exercicios}
  \item $2^3 + 3^2$
  \item $5^2 - 4^2$
  \item $10 - 2 \times 2^2$
  \item $(10 - 2 \times 2)^2$
  \item $\sqrt{64} + \sqrt[3]{8}$
  \item $3 \times \sqrt{36} - 2^4$
  \item $(-2)^4 - 2^4 + \sqrt{81}$
  \item $\sqrt{5^2 - 3^2}$
  \item $2^5 : 2^3 + \sqrt[3]{125}$
  \item $\left(\sqrt{49} - 2\right)^3 : 5^2$
\end{exercicios}

\secao{Desafios}{Resolve os problemas.}
\begin{exercicios}
  \item Qual é o algarismo das unidades de $2^{10}$?
  \item Determina $n$ de modo que $3^n = 243$.
  \item Determina $n$ de modo que $2^n \times 2^3 = 2^{10}$.
  \item Escreve 64 como potência de base 2, de base 4 e de base 8.
  \item Uma bactéria divide-se em duas a cada hora. Começando com uma bactéria, quantas existem ao fim de 8 horas? Escreve o resultado como potência.
  \item Um terreno quadrado tem 225\,m$^2$. Quantos metros de rede são necessários para o vedar?
  \item Um cubo de aresta 3\,cm é formado por cubinhos de 1\,cm de aresta. Quantos cubinhos são? E se a aresta for 5\,cm?
  \item Com $a = 16$ e $b = 9$, verifica que $\sqrt{a \times b} = \sqrt a \times \sqrt b$, mas $\sqrt{a + b} \neq \sqrt a + \sqrt b$.
\end{exercicios}

\end{document}
